\clearpage
\VerdierPhysicalPageAnchor{0258}%
\begingroup
\small
\setlength{\parindent}{0pt}%
\setlength{\parskip}{0pt}%
\setlength{\emergencystretch}{2em}%
\newcommand{\VerdierTermRow}[3]{%
  \VerdierSemanticAnchor{#1}%
  \par\noindent\hangindent=1.5em\hangafter=1
  #2。\enspace\textup{#3。}\par\smallskip}

\VerdierSemanticAnchor{idxt-title}
\vspace*{2.5em}
\begin{center}
  {\Large\bfseries 用語索引}
\end{center}
\vspace{2em}

\VerdierTermRow{idxt-p258-r01}{スペクトル系列の収束先}{第 II 章、4.4.3}
\VerdierTermRow{idxt-p258-r02}{シフト自己同型}{第 II 章、1.1.2}
\VerdierTermRow{idxt-p258-r03}{$\beta$-射（捩れ $G$-関手の間の）}{第 I 章、1.4.6}
\VerdierTermRow{idxt-p258-r04}{分解可能な圏}{第 II 章、1.2.8}
\VerdierTermRow{idxt-p258-r05}{アーベル圏の導来圏}{第 III 章、1.2.2}
\VerdierTermRow{idxt-p258-r06}{$C$ の複体のホモトピー圏}{第 I 章、2.5.7}
\VerdierTermRow{idxt-p258-r07}{$n$-重複複体のうち $\varphi$-総和可能なものの圏}{第 I 章、2.1.7}
\VerdierTermRow{idxt-p258-r08}{加法的次数付き圏}{第 I 章、1.1.3}
\VerdierTermRow{idxt-p258-r09}{$G$ 型次数付き圏}{第 I 章、1.1.1}
\VerdierTermRow{idxt-p258-r10}{乗法的系に関して局所化された圏}{第 II 章、2.2.5}
\VerdierTermRow{idxt-p258-r11}{$G$-圏の反対圏}{第 I 章、1.3.9}
\VerdierTermRow{idxt-p258-r12}{前三角圏}{第 II 章、2.2.12}
\VerdierTermRow{idxt-p258-r13}{三角圏の三角部分圏による商圏}{第 II 章、2.2.10}
\VerdierTermRow{idxt-p258-r14}{三角圏}{第 II 章、1.1.1}
\VerdierTermRow{idxt-p258-r15}{反対三角圏}{第 II 章、1.1.7}
\VerdierTermRow{idxt-p258-r16}{$f$ のホモトピー類}{第 I 章、2.5.7}
\VerdierTermRow{idxt-p258-r17}{基底に関する Koszul 類}{第 I 章、1.6.3}
\VerdierTermRow{idxt-p258-r18}{順序付き基底に関する Koszul 余サイクル}{第 I 章、1.6.4}
\VerdierTermRow{idxt-p258-r19}{非輪状複体}{第 III 章、1.1.3}
\endgroup
